% Self-contained report section.  No claim of a completed 245-state machine.
\section{Shortening the return scanner: an exact obstruction and a source target}
\label{sec:scanner-eight}

The return scanner offers a concrete connection between source-tree
reassociation and the physical state count.  The current 246-state
machine recognizes nine consecutive zeros while moving left.  Replacing
this with an eight-zero recognizer removes one working state, but the
change requires a stronger condition on the words that the interpreter
scans.  The local scanner theorem and the exact obstruction below are
proved; they do not assert that this isolated 245-state table implements
the complete program.

\paragraph{The literal patch.}
Use the original state labels of
\texttt{zfc\_prefix246\_old.tm}.  Delete both rows of state 74 and replace
the single incoming row
\[
 (58,0)\longmapsto(0,L,74)
 \qquad\text{by}\qquad
 (58,0)\longmapsto(0,L,84).
\]
All remaining rows are identical to the 246-state table.  State 74 has no
other incoming transition, so the result is a total binary transition
table with 245 working states and the same halt state 254.  States 4 and
6 are retained because the register-marking contracts also use them.

Let
\[
 (s_0,s_1,\ldots,s_7)=(4,6,9,13,19,29,41,58).
\]
The patched rows are exactly
\[
 \delta(s_i,1)=(1,L,s_0),\qquad
 \delta(s_i,0)=
 \begin{cases}
  (0,L,s_{i+1}),&0\leq i<7,\\
  (0,L,84),&i=7.
 \end{cases}
\]
State 84 has the unchanged rows
\[
 \delta(84,0)=(0,R,84),\qquad
 \delta(84,1)=(1,R,96).
\]
Thus $s_i$ records that $i$ consecutive zeros have already been read
while moving left.  Every displayed transition preserves its tape cell.

\paragraph{Universal local scanner lemma.}
Fix integers $m\leq x$.  Suppose that the tape cell at $m$ is one, all
cells strictly left of $m$ are zero, and the finite word in the interval
$[m,x]$ contains no occurrence of $0^8$.  From state 4 with head at $x$,
the patched table reaches state 96 with head at $m+1$, leaves the entire
tape unchanged, and takes exactly
\[
 (x-m+1)+18=x-m+19
\]
literal transitions.  The contents strictly right of $x$ are arbitrary.

\paragraph{Proof.}
Induct on the number of cells already scanned in $[m,x]$.  The state is
$s_i$, where $i$ is the length of the consecutive zero suffix of the
word read in the direction of travel.  A one resets this counter to
zero; a zero increases it by one.  The forbidden-word hypothesis keeps
the counter below eight until the marker at $m$ has been read.  After
exactly $x-m+1$ transitions, the head is at $m-1$ in state $s_0$.
Eight exterior zeros then place the head at $m-9$ in state 84.  The
rightward self-loop crosses the nine zeros from $m-9$ through $m-1$,
and the marker-one row moves once more to $(96,m+1)$.  The final ten
transitions give the stated total.  Every row preserves its symbol and
no cell to the right of the entry head is read.\hfill$\square$

The same argument applies when entry is in state $s_i$ rather than
$s_0$, provided the interval word followed by $i$ virtual zeros contains
no $0^8$.  In particular, state 6 means that one zero has already been
counted.  Its valid hypothesis is the absence of $0^8$ in the interval
word followed by one virtual right-hand zero; absence of $0^8$ in the
interval word alone is insufficient.  The marker resets the counter,
so the transition count after consuming the interval is still the
same 18.  Relative to the original nine-zero scanner, a successful
return from state 4 saves exactly two transitions.

\paragraph{The register contracts are preserved.}
The four one-block induction proofs, the free-word identity, all three
increment/decrement contracts, and all nine register-marking proofs
replay against the patched literal table.  Their derivations do not use
state 74 or the changed row $(58,0)$.  They continue to quantify over
arbitrary natural register values, every finite prefix and suffix of
blocks, arbitrary left context, and a nonempty suffix after the selected
primary register.  This establishes the register premises needed by the
affine interpreter checker independently of the local scanner lemma.

\paragraph{The baseline audit identifies the precise failure.}
Apply the existing affine symbolic instruction checker to every leaf
of the original pinned \texttt{Prog254.v}, using the patched table.
Its frame is
\[
 {}^\infty0\;1\;P\;0^{\,d-|P|}\;1\;B(v)\;0^\infty,
 \qquad q=96,\quad d=40+n,\quad n\geq0,
\]
where the head is immediately after the left marker, $v$ consists of
nine arbitrary natural-valued primary registers followed by an
arbitrary nonempty finite spare list, and $B(a)=01^{a+1}$.
The lower bound 40 concerns the whole head-to-header distance $d$;
it is not a lower bound on the trailing zero separator alone.

Of the 8,391 guarded instruction cases, 8,355 have complete successful
symbolic derivations for all admitted frames.  The breakdown is
\[
\begin{array}{c|r|r}
 \text{case}&\text{successful}&\text{failed}\\\hline
 \mathrm{CI}&2601&0\\
 \mathrm{CD},\ \text{zero}&1913&18\\
 \mathrm{CD},\ \text{positive}&1913&18\\
 \mathrm{CB}&128&0\\
 \mathrm{CC}&1799&0\\
 \mathrm{CH}&1&0\\\hline
 \text{total}&8355&36
\end{array}
\]
Every failed case belongs to one of exactly 18 decrement leaves whose
path ends with seven zeros.  There are no failures in any other
instruction class.  These are successful or failed symbolic proof
derivations, rather than a finite sample of concrete register values.

The ordinary decrement phase protocol explains this obstruction.
Preparation changes the two phase cells from $00$ to $11$ and fills
the remaining separator with ones.  The register-selection step
changes the first phase one to zero, producing $01$.  During the
subsequent return, the first phase zero is adjacent to the last program
counter bit.  Therefore a path ending in $0^7$ produces a scanned
$0^8$ run.  The shortened scanner turns around there before reaching
the left marker.  Merely observing that the source path itself has no
$0^8$ misses this extra phase zero.

\paragraph{How the new source graph addresses the obstruction.}
The separately certified source graph
\texttt{rotation\_zero\_candidate\_24.json} has maximum path depth 36,
202 distinct shared pair subtrees, and no occurrence of $0^7$ in any
instruction path.  Appending one phase zero to such a path gives a
maximum zero run of seven, which satisfies the eight-zero local lemma.
This is a useful source-side target for the compiler.  The source bound
does not, by itself, prove that every compiler alias or intermediate
phase meets the scanner precondition.  The literal compilation must
also be checked, including its intermediate phases.

\paragraph{A complete universal runtime check at 250 states.}
A repaired compilation of that source graph has 251 working states.
It explicitly splits three conflicting shared decoder roles and
repairs the affected break-instruction return chains.  Applying the
single scanner patch above gives
\texttt{compiled24\_scanner250.tm}, a 250-state literal table.
All four block inductions, the word identity, all three arithmetic
macros, and all nine marking contracts pass against this exact table.
The full affine audit also succeeds at every one of the 8,391 guarded
instruction cases:
\[
\begin{array}{c|rrrrrr}
 \text{case}&\mathrm{CI}&\mathrm{CD}/0&\mathrm{CD}/+&
              \mathrm{CB}&\mathrm{CC}&\mathrm{CH}\\\hline
 \text{proved}&2601&1931&1931&128&1799&1.
\end{array}
\]
Each successful case quantifies over every $n\geq0$, every
natural-valued register vector satisfying its named decrement guard,
and every nonempty finite spare suffix.  The initial frame uses
$d=40+n$, even though this source has maximum depth 36.  The audit
therefore establishes that the existing lower bound $d\geq40$ is
sufficient; increasing it to 42 is unnecessary.  Successful
nonhalting cases reach the entire expected successor frame, including
the exact register effect, spare suffix, program counter, head
position, and header distance.  The halt case reaches state 254.
Every case has positive literal progress; the smallest literal-step
count in a symbolic derivation is 13.  The certificate is
\texttt{compiled24\_scanner250\_audit.json}.

This proves the full guarded runtime simulation of the rotated source
by that table from admitted instruction frames.  The independent
source-equivalence certificate consequently transfers the original
register program's behavior at corresponding source entries.
Blank initialization is now certified separately by
\texttt{agent\_states/check\_blank\_prefix.py} and
\texttt{agent\_states/blank\_prefix\_results.json}.  From state 95
on the all-zero tape, this exact 250-state table reaches the reference
safe root after 407,983 transitions.  The endpoint is state 96 with
head $-5$, marker $-6$, header 176, empty PC, physical primary vector
$(0,0,0,3,0,0,3,0,0)$, and 167 spare blocks, each of value one.
The complete tape has 351 one-cells, and every other cell is zero.
The fixed-prefix check verifies the whole tape and that cells outside
the explicit finite window were never visited.  This supplies the initialization component in
addition to the separately proved runtime and source correspondence.
The table still has 250 working states, above the requested bound.

The supplied \texttt{check\_short\_scanner.py} checks the exact scanner
rows, replays the register inductions, expands an optional source graph
with structural checks, and checks every instruction frame.  Its
baseline certificate and compact summary record the failed cases
explicitly, and the same checker produces the successful 250-state
certificate.  The isolated 245-state baseline scanner patch remains
a rejected whole-program candidate.  The validated 250-state runtime
shows that source reassociation can discharge the scanner's universal
precondition, while also making the cost of repairing physical decoder
aliases explicit.  A below-246-state solution still requires a smaller
complete table and its blank-tape initialization argument.
